% !PRINTED 129
\weberpage{129}
%<<CONT
\weq{17}{
E\!\left(\begin{matrix}\alpha+\beta, & \beta \\ \gamma+\delta, & \delta\end{matrix}, \omega\right)
= \left(\frac{\pm\beta}{\pm(\alpha+\beta)}\right)
i^{-\frac{1\mp(\alpha+\beta)}{2}}
e^{\frac{\pi i}{12}\left\{(\alpha+\beta)(\gamma+\delta-\beta)-[(\alpha+\beta)^2-1]\beta\delta\right\}}
\sqrt{\pm(\alpha+\beta+\beta\omega)}.
}
Nun ist, wenn die unteren Zeichen gelten, $\beta$ negativ, und
wenn also
\[
-(\alpha+\beta+\beta\omega) = r\,e^{i\varphi}
\]
\[
i(\alpha+\beta+\beta\omega) = r\,e^{i\left(\varphi-\frac{\pi}{2}\right)}
\]
gesetzt wird, so liegt $\dfrac{\varphi}{2}$ und $\dfrac{\varphi}{2}-\dfrac{\pi}{4}$ zwischen $-\dfrac{\pi}{2}$
und $+\dfrac{\pi}{2}$, und daraus folgt, daß wir zu setzen haben:
\[
\sqrt{-(\alpha+\beta+\beta\omega)} = i\sqrt{\alpha+\beta+\beta\omega},
\]
ferner nach dem Reziprozitätsgesetz der quadratischen Reste:
\[
\left(\frac{\beta}{\alpha+\beta}\right) = (-1)^{\frac{\beta(\alpha+1)}{4}}\left(\frac{\beta}{\alpha}\right),
\]
\[
\left(\frac{-\beta}{-(\alpha+\beta)}\right) = (-1)^{\frac{\alpha-1}{2}}(-1)^{\frac{\beta(\alpha-1)}{4}}\left(\frac{\beta}{\alpha}\right)\,^{1}),
\]
und daraus folgt die Übereinstimmung der beiden Ausdrücke
(16), (17) und mithin die Richtigkeit der Formel (13) nach der
Kongruenz
\[
\beta(3\alpha-2\gamma+\beta-\delta+\beta^2\delta+2\alpha\beta\delta) \equiv 0\ (\mathrm{mod}\ 24),
\]
die sich, da $\beta$ gerade vorausgesetzt ist, aus $\alpha\delta-\beta\gamma=1$ ergibt.

Somit sind also die Formeln (15) als richtig erwiesen.

\noindent\rule{4cm}{0.4pt}

\noindent $^{1})$ Nach dem Reziprozitätsgesetz ist, wenn $\beta=\pm2^\lambda\beta'$ gesetzt und $\beta'$
ungerade und positiv angenommen wird, wenn $\alpha+\beta$ positiv ist
\[
\left(\frac{\beta}{\alpha+\beta}\right) = \left(\frac{\pm2^\lambda}{\alpha+\beta}\right)\left(\frac{\alpha}{\beta'}\right)
(-1)^{\frac{(\alpha+\beta-1)(\beta'-1)}{4}},
\]
\[
\left(\frac{\beta}{\alpha}\right) = \left(\frac{\pm2^\lambda}{\alpha}\right)\left(\frac{\alpha}{\beta'}\right)
(-1)^{\frac{(\alpha-1)(\beta'-1)}{4}}.
\]
Ist $\lambda\geqq2$, so sind diese beiden Werte einander gleich, ist $\lambda=1$, so
unterscheiden sie sich durch den Faktor $(-1)^{\frac{\alpha+1}{2}}$, in Übereinstimmung
mit der ersten der obigen Formeln. Die Richtigkeit der zweiten Formel
ergibt sich, wenn $\alpha+\beta$ negativ ist, aus
\[
\left(\frac{-\beta}{-(\alpha+\beta)}\right) = \left(\frac{2^\lambda}{-(\alpha+\beta)}\right)\left(\frac{\alpha}{\beta'}\right)
(-1)^{\frac{(\alpha+\beta-1)(\beta'-1)}{4}},
\]
\[
\left(\frac{\beta}{\alpha}\right) = \left(\frac{2^\lambda}{\alpha}\right)\left(\frac{\alpha}{\beta'}\right)(-1)^{\frac{\alpha-1}{2}}
(-1)^{\frac{(\alpha-1)(\beta'-1)}{4}}.
\]
